\section{The positive anchor and the continued filling mark}
\label{ng:anchor}

We first identify the family at $M=1$ with the original positive
geometric family, including the endpoint of its filling path.
This step does not follow just from agreement at the complete point.

Put $Q=e^{\ii t}$, $P=e^{\ii u}$ with $|t|,|u|\le10^{-2}$,
and write $y=\tan(t/2)$, $z=\tan(u/2)$ and
$k=(1+y^2+z^2)^{-1/2}$, with $k=1$ at zero.
The angular chart of \cite[Section~3, equations~(3.3)--(3.5)]{Circle}
has
\begin{equation}\label{ng:angular-match}
 E=\frac{1+\ii k}{1-\ii k},\qquad
 R=\frac{1+z^2}{2+y^2+z^2},\qquad
 H_0=\frac{1-ky}{1+ky},\quad H_1=\frac{1-kz}{1+kz}.
\end{equation}
At $m=1$, equation \eqref{ng:N16} becomes
\[
 X+X^{-1}=\frac{4-\beta\gamma}{\beta+\gamma+4}
       =\frac{2(y^2+z^2)}{2+y^2+z^2}.
\]
Thus $X=E_0=E$ on the branch through $\ii$.
The first apex is
$(E+Q)/(1+Q)=(1+E)/(1+H_0)$.
Writing $Z_1=(1+E)/(1+H_1)$, the second is
$P(1-E)/(P+1)=RZ_1/(Z_1-(1-R))$.
These are exactly the ordered apices of the angular block.
The four matrices \eqref{ng:N3} specialize projectively to its
four white maps. For example $W_2$ differs from the representative
in \cite[equation~(3.8)]{Circle} only by a scalar.
The pure shaded maps agree because they have the same three ordered
images. Equation \eqref{ng:S14} identifies the odd collar for all
nearby $E,Q$, so it also agrees along this angular family.
The entire marked representation, rather than just its peripheral
trace, is therefore the same one.

For clarity, these formulas are regular throughout the stated angular
bidisk. The elementary Taylor bounds give
$|Q-1|,|P-1|<11/1000$ and
$|\beta-2|,|\gamma-2|<1/9000$.
Indeed $e^r-1\le r/(1-r)$ and
$2(\cosh r-1)\le r^2/(1-r^2)$ for $r=1/100$.
Hence $|X+X^{-1}|<1/16000$; the root through $\ii$ has
$|X-\ii|<1/1000$.
For the last assertion write
$X=A_0/2+\ii\sqrt{1-A_0^2/4}$ with $A_0=X+X^{-1}$,
using the square-root branch at one.
At $m=1$,
$R=(X+X^{-1}+\beta)/(\beta+2)$, so $|R-1/2|<1/1000$.
The apices in \eqref{ng:N5} differ from their complete values by
less than $1/50$. The finite defining vertices therefore stay
distinct, and all displayed denominators stay nonzero.
Moreover
\[
 G|_{m=1}=PQ(\beta+\gamma+4)
              \{E_0^2-(X+X^{-1})E_0+1\}
\]
has a simple root on this bidisk. The change \eqref{ng:N12}
has derivative one in $E_0$ at $m=1$.
Thus the native charts extend holomorphically in $M$ around every
point of the angular family. For a fixed finite tree, compactness of
its charge path supplies finitely many overlapping such charts.
No unspecified complete-point radius is substituted for this path.

Let $\Psi(\boldsymbol\tau)=\sum_{vw}\psi(\tau_v,\tau_w)$ be the
classical angular potential of \cite[equation~(1.3)]{Circle};
its normalization and its derivatives are
\[
 \psi(0,0)=v_{\rm oct}/\pi,\qquad
 \psi_t(t,u)=-\frac1\pi
       \operatorname{arsinh}\bigl(\tan(t/2)\cos(u/2)\bigr).
\]
The actual auxiliary increments in that paper give
\begin{equation}\label{ng:anchor-logs}
 U_v=-2\pi\Psi_{\tau_v}-\ii\epsilon_v\tau_v,\qquad
 V_v=2\ii\tau_v,\qquad Q_v=e^{\ii\tau_v}.
\end{equation}
This is the specialization of \eqref{ng:T3}, with the same
curves and logarithms, since the local maps just matched.
The charge path satisfies
\begin{equation}\label{ng:charge-path}
 \Psi_{\tau_v}+\frac{\ii n_v}{2\pi}\tau_v=\ii c
 \quad\Longleftrightarrow\quad
 U_v+p_vV_v=-2\pi\ii c,\qquad 0\le c\le1.
\end{equation}
The logarithms are continued from $c=0$, not reset to principal
values. The constant on the right records one primitive winding
at $c=1$.

Under either geometric criterion in Theorem~\ref{ng:filled-germ},
\cite[Theorems~1.1 and~7.1]{Circle} constructs this positive
path, identifies the full log-link periods, and proves smooth
completion at $c=1$ on the original filled exterior.
It also proves $\operatorname{Im}\tau_v(c)<0$ for $c>0$.
Consequently $|Q_v|>1$ there, so every filled longitude has
distinct eigenvalues. Intermediate $0<c<1$ has cone angle
$2\pi c$ and generally does not descend to the smooth filled group.
Only the endpoint is used below.

\section{Local rigidity on the retained augmented triangulation}
\label{ng:rigidity}

We give the local rigidity argument with hypotheses that can be checked
on another original positive augmented triangulation as well.
There is no assertion that an arbitrary hyperbolic manifold has the
particular positive triangulation required here.

\begin{proposition}[Marked positive gluing]\label{ng:positive-germ}
Let an oriented ideal triangulation of a specified manifold have $t$
tetrahedra, all with positive shapes, and $c$ torus ends. Retain all
actual face pairings, peripheral paths, developed flags and continuous
edge-log branches. The ends may be incomplete. Suppose the edge
equations hold and all face triples are distinct.
Choose one nonzero primitive peripheral class on each torus.
Its $c$ similarity logarithms are local holomorphic coordinates on
the decorated representation germ modulo simultaneous projective
conjugation. The shape and decorated-representation descriptions are
inverse also over local Artin algebras with this residue point.
\end{proposition}

\subsection{The real row space}
For each positively ordered tetrahedron use
\begin{equation}\label{ng:shape-coordinates}
 u_j=\log z_j,\qquad v_j=-\log(1-z_j),\qquad
 T_j=\frac{\dd v_j}{\dd u_j}=\frac{z_j}{1-z_j},
 \qquad \operatorname{Im}T_j>0 .
\end{equation}
Logarithms continue from the given solution. The original edge and
peripheral equations are affine linear in these logs. Their constants
need not be small.

Here is the rank calculation. The number of global edges is $t$:
coning off the $c$ torus boundaries gives Euler characteristic $c$,
whereas its cell count is $c-e+2t-t$. A linear dependence among edge
rows assigns numbers $\lambda_{ij}$ to the six edges of each
tetrahedron with the three opposite-pair sums equal.
Those equations are equivalent over $\RR$ to
$\lambda_{ij}=s_i+s_j$. For example
$s_0=(\lambda_{01}+\lambda_{02}-\lambda_{12})/2$ and the
other $s_i$ follow; the opposite-sum conditions give the remaining
three edges. The three edge values on a face determine its three
$s_i$ uniquely, so the values descend to one per ideal vertex.
Conversely every such vertex assignment is a dependence.
The assignment is injective, since the three zero endpoint sums
on one triangle force all three vertex values to vanish.
The edge-row space $\mathcal E$ therefore has dimension $t-c$.

We use the actual normal-path identity proved in
\cite[Section~7.1, equations~(7.1)--(7.2)]{Marked}.
It identifies each path row with its developed similarity logarithm
and proves that the row pairing is twice oriented intersection,
up to its fixed coordinate sign. Briefly, the local outgoing-flow
formula is a sum of two-coordinate determinants plus one skew form
on each face. Paired face forms cancel under the orientation-reversing
face identifications. A small normal circle about an edge endpoint
is that edge's incidence row and is null-homologous on the cusp.
Thus $\mathcal E$ is isotropic and orthogonal to every peripheral
row. The $2c$ rows of primitive peripheral bases are independent
in $\mathcal E^\perp/\mathcal E$, because their pairings are
nonzero on each torus and zero between different tori.
Its dimension is $2c$, so those rows form a basis.

The cited calculation uses $(u,v_{\rm old})=(\log z,\log(1-z^{-1}))$.
Here $v_{\rm old}=-u-v+\text{constant}$; the real coordinate
change has determinant $-1$. It reverses the pairing sign but
preserves all ranks and isotropy. Thus the same conclusions hold
for \eqref{ng:shape-coordinates}. No longitude divisibility or
special knot homology assertion is used.

The chosen one-per-torus classes are independent modulo
$\mathcal E$ and pair to zero. Append them to $t-c$ independent
edge rows. The resulting real $t$-row matrix $(A\ B)$ satisfies
\begin{equation}\label{ng:real-rows}
 \operatorname{rank}(A\ B)=t,\qquad
 AB^{\mathsf T}=BA^{\mathsf T}.
\end{equation}
Its kernel is spanned by $(-B^{\mathsf T}y,A^{\mathsf T}y)$ and
is a real Lagrangian space for
$\Omega((a,b),(a',b'))=a\cdot b'-b\cdot a'$.

\subsection{Positive polarization and the actual representation germ}
The selected logarithmic Jacobian is
$J=A+B\diag(T_j)$.
If $Jx=0$, then
$w=(x,\diag(T_j)x)$ and $\overline w$ belong to the
complexified real Lagrangian kernel. But
\begin{equation}\label{ng:polarization}
 \frac{\ii}{2}\Omega(w,\overline w)
              =\sum_j\operatorname{Im}T_j\,|x_j|^2 .
\end{equation}
Its left side is zero, and its right side is positive for $x\ne0$.
Thus $J$ is invertible. Dependent edge rows have compatible affine
constants because the given point satisfies all original equations.
The holomorphic implicit-function theorem proves the asserted
coordinate statement for the shape germ.

This is the actual decorated representation germ. Normalize one
ordered developed triple. The shapes develop successive tetrahedra
across their paired faces by the unique map of three distinct
vertices. An edge return fixes its two distinct endpoints and
has multiplier one by the actual edge equation; hence it is
projectively identity. The dual two-skeleton of the original
triangulation presents its fundamental group. The constructed
face maps therefore give a representation of that group, with
the prescribed peripheral fixed flags and path logarithms.

Conversely a nearby representation $\rho$ with selected peripheral
flags $\xi_v$ develops every ideal vertex $gP_v$ as
$\rho(g)\xi_v$, where $P_v$ is its peripheral subgroup.
This is well defined because $P_v$ fixes $\xi_v$.
The distinct face triples and tetrahedral quadruples remain
distinct nearby. Their cross-ratios recover the shapes; the
face transformations agree on three points, and edge products
telescope. The two constructions are inverse after the same
triple normalization. All denominators are units at the given
point, so they remain inverse over local Artin algebras.
Here an Artin algebra is a finite-dimensional local complex
algebra with nilpotent maximal ideal; it records infinitesimal
as well as ordinary nearby deformations.

At a parabolic peripheral matrix it is essential to retain the
exact eigenvalue character. For
$\left(\begin{smallmatrix}1&1\\0&1\end{smallmatrix}\right)$,
the invariant line generated by $(1,\eta)$ can move when
$\eta^2=0$, but its eigenvalue becomes $1+\eta$.
Fixing the exact eigenvalue one forces $\eta=0$.
In general a unit off-diagonal entry gives the selected eigenvector
$(\mu_{12},\lambda-\mu_{11})$ and the same unique free eigenline.
The similarity logs in the proposition retain that character,
not only the unmarked trace. A separately prescribed determinant-one
lift keeps its central signs on its local sign cover.
This proves Proposition~\ref{ng:positive-germ}.

\section{Filling words and the one-dimensional geometric germ}
\label{ng:filling-germ-proof}

At the augmented complete structure choose the main meridian and
the auxiliary longitudes in Proposition~\ref{ng:positive-germ}.
Their logarithms are $-2\log M$ and $2\log Q_v$ in our selected
flags. The explicitly constructed family has these very
coordinates. It therefore exhausts the nearby decorated
augmented representation germ. This completes the local assertion
of Theorem~\ref{ng:representation}. The continuation and positive
anchor in Section~\ref{ng:anchor} prove its endpoint assertion.

Now let $\rho_*$ be a positive primitive filled endpoint as in
Theorem~\ref{ng:filled-germ}. Use the original augmented
triangulation at this endpoint. Its auxiliary ends are incomplete;
the proof of Proposition~\ref{ng:positive-germ} did not require
their completeness or the faithfulness of the augmented representation.
Choose the main meridian and the original slopes
$\sigma_v=\mu_v+p_v\lambda_v$ as its peripheral coordinates.
Each is a primitive nonzero class on a different torus.
Their log rows are real and satisfy exactly the hypotheses of
the proposition. Fixing the $R$ auxiliary slope logs leaves a
smooth one-dimensional shape and decorated representation germ.

We must still identify this condition with the filled group.
A slope word commutes projectively with the auxiliary longitude.
The latter has distinct eigenvalues by hypothesis. Choose a local
representative of the slope whose residue is identity. If its
commutator with the longitude has scalar $c_0$, determinants give
$c_0^2=1$, while $c_0\equiv1$ at the residue point.
Because $c_0+1$ is a unit, $c_0=1$, also over local Artin algebras.
The longitude eigenspaces split the rank-two module because their
eigenvalue difference is a unit. Strict commutation kills both
off-diagonal entries of the slope in that basis.
Equation \eqref{ng:A2} gives eigenvalue ratio one, so the slope
matrix is scalar and hence projectively identity.
This normalization does not reset the signs of a separately
retained determinant-one group lift.

Conversely, a nearby representation of the filled group pulls
back to one whose slopes are projectively identity.
On the continued branch at $\rho_*$ their logarithms are precisely
$-2\pi\ii$, not another multiple of $2\pi\ii$.
The auxiliary eigenflags are uniquely selected by the longitude
branches. At the main cusp retain the exact eigenvalue $M^{-1}$;
the formula for its eigenvector in the preceding proof remains
regular at its noncentral parabolic basepoint.
Thus filled representations and the fixed-slope decorated shape
germ are the same local object, modulo simultaneous conjugation.
Its coordinate is $\log M$, since the main similarity logarithm
on that line is $-2\log M$, a nonzero scalar multiple.

For the tree endpoints, Section~\ref{ng:anchor} supplied the
original smooth metric holonomy and the nonunit auxiliary
longitudes. This proves Theorem~\ref{ng:filled-germ}, without
assuming rigidity from a new positive triangulation of the filled
manifold.

\begin{corollary}[Identification of an independently constructed branch]
\label{ng:branch-identification}
Suppose another holomorphic construction reconstructs the same
original augmented representation at this positive endpoint,
keeps its main eigenflag and primitive logarithmic slope words,
and satisfies \eqref{ng:A2} nearby with the same variable $M$.
Its local representation branch is the geometric meridian germ
of Theorem~\ref{ng:filled-germ}. This holds also for infinitesimal
branches with the exact marked characters retained.
\end{corollary}
\begin{proof}
Both branches lie in the same decorated filled germ, and $\log M$
is a local coordinate on it, including over local Artin algebras.
\end{proof}

In a stationary-source application, nondegeneracy of that source
critical point is one way to supply the independent holomorphic
branch in the corollary. It is not a premise of the geometric
construction or a consequence of a volume derivative. The
preferred main longitude, its framing monomial and a quantum
annihilator remain separate identifications.
